\section{Orthogonal splitting of positive semidefinite matrices}
\label{sec:splitting}

The next lemma turns small pairwise overlaps of positive semidefinite
matrices into an orthogonal decomposition of their common ambient space.
Each summand of the decomposition is assigned one matrix, and the lemma
controls the total squared norm of the other matrices compressed to that
summand.  The bound depends on the number of matrices, but not on their
dimension.  We will apply it to Gram matrices of Fourier coefficients.

\begin{lemma}[Orthogonal splitting]\label{lem:splitting}
Let $A_1,\ldots,A_L$ be real symmetric positive semidefinite matrices
on a finite-dimensional Euclidean space.  There are mutually orthogonal
projections $P_1,\ldots,P_L$ such that $\sum_{s=1}^L P_s=I$ and
\begin{equation}\label{eq:orthogonal-splitting}
  \sum_{z=1}^L\sum_{s\ne z}\normF{P_sA_zP_s}^{2}
  \le C_L\sum_{z\ne z'}\tr(A_zA_{z'}),
  \qquad C_L=2^{L-1}.
\end{equation}
The sum on the right is over ordered pairs of distinct indices.
Zero projections are allowed.
\end{lemma}

\begin{proof}
We first record the two-matrix estimate that gives the decomposition.
For positive semidefinite $A,B$, let $P$ be the spectral projection of
$A-B$ onto its nonnegative eigenspaces, and let $Q=I-P$.
This is the same spectral projection used in binary quantum
discrimination~\cite{helstrom}; here we need a bound on squared
Frobenius norms of the compressed matrices.
In the block decomposition associated with $P$ and $Q$, write
$A=(A_{uv})_{u,v=1}^2$ and $B=(B_{uv})_{u,v=1}^2$.
Since $P$ commutes with $A-B$, we have
$A_{12}=B_{12}$, $A_{11}\succeq B_{11}$, and
$B_{22}\succeq A_{22}$.  Thus
\begin{equation}\label{eq:two-matrix-splitting}
\begin{split}
  \tr(AB)
  &=\tr(A_{11}B_{11})+\tr(A_{22}B_{22})
       +2\normF{A_{12}}^2\\
  &\ge \normF{PBP}^2+\normF{QAQ}^2.
\end{split}
\end{equation}
Here we used the nonnegativity of the trace of a product of two
positive semidefinite matrices.

We now induct on $L$.  The assertion for $L=1$ holds with $P_1=I$.
For $L\ge2$, apply the preceding construction to
$A=A_1$ and $B=\sum_{j=2}^L A_j$, and put $P_1=P$.
Apply the induction hypothesis to the compressions $QA_jQ$, $2\le j\le L$,
on the range of $Q$.  Extend the resulting projections by zero on the
range of $P$, obtaining $P_2,\ldots,P_L$ on the original space.

The terms on the left of \eqref{eq:orthogonal-splitting} involving
either $P_1$ or $A_1$ have sum at most
\begin{equation}\label{eq:splitting-first-label}
  \sum_{j=2}^L\normF{PA_jP}^2
    +\sum_{s=2}^L\normF{P_sA_1P_s}^2
  \le \normF{PBP}^2+\normF{QA_1Q}^2
  \le \tr(A_1B).
\end{equation}
The first inequality uses nonnegative pairwise trace products in the
first sum and the orthogonal block decomposition of $QA_1Q$ in the
second.  To bound the overlaps introduced by compression, observe that
for distinct $j,j'\ge2$,
\begin{equation}\label{eq:compressed-overlap}
\begin{split}
  \tr(A_jQA_{j'}Q)
  &=\normF{A_j^{1/2}QA_{j'}^{1/2}}^2\\
  &\le 2\tr(A_jA_{j'})+2\tr(A_jPA_{j'}P).
\end{split}
\end{equation}
This is the squared triangle inequality applied to $Q=I-P$.
Moreover,
\[
  \sum_{\substack{j,j'\ge2\\j\ne j'}}\tr(A_jPA_{j'}P)
  \le \normF{PBP}^2
  \le \tr(A_1B).
\]
Combining these estimates with the induction hypothesis bounds the
left side of \eqref{eq:orthogonal-splitting} by
\[
  (1+2C_{L-1})\tr(A_1B)
  +2C_{L-1}\sum_{\substack{j,j'\ge2\\j\ne j'}}\tr(A_jA_{j'}).
\]
The full ordered overlap sum is
$2\tr(A_1B)+\sum_{j\ne j',\,j,j'\ge2}\tr(A_jA_{j'})$.
Since $C_{L-1}\ge1$, the preceding bound is at most
$2C_{L-1}$ times this sum.  Taking $C_L=2C_{L-1}$ completes the induction.
\end{proof}
